\scanpage{173}
% L1–8
$\pi:F\mathbin{\text{---}}E$ telle que
\[
\epsilon\pi=\lambda\,\mathrm{id}_F,
\]
cela signifie que l'extension
$ {}^{\lambda}E$ de $F$ par $E$, déduite de $E$ par
image inverse par l'hom. $\lambda\,\mathrm{id}_F$ \unsure{de|dans} $F$,
ou par image directe par $\lambda\,\mathrm{id}_E$ de $E$,
est munie d'un scindage.
On dit que $E$ est $\lambda$-scindée.

% L9–14
Si p.ex. $F=k$ Cas \uline{d'}\uline{un} espace
affine $A=\epsilon^{-1}(1)$ sous $E^0$), la
donnée de $\pi$ \inserted{$:k\unsure{\mathbin{\text{---}}|\longrightarrow} E$} revient à la donnée
d'un élément $t$ de $E$, et la
condition envisagée c'est que
\[
\epsilon(t)=\lambda .
\]

% L15–21
Cela signifie aussi que le torseur
\[
\epsilon^{-1}(\deleted{\unsure{1|A}}\lambda)={}^{\lambda}A
\]
déduit de $A$ par l'homothétie $\lambda\,\mathrm{id}_{A^0=E^0}$
sur le module des translations, est
muni d'une section \inserted{``origine''}. Donc
l'application
\[
\begin{aligned}
x&\longmapsto\lambda x && A\mathbin{\text{---}}{}^\lambda A\\
&\text{compatible avec} && A^0\mathrel{\mathop{\text{---}}\limits^{\lambda\,\mathrm{id}}} A^0
\end{aligned}
\]
% archive sheet number: 172
